\begin{abstract}

Whether a language model ``respects logic'' has no answer until three
things are fixed: which logical identity, which behavioral identity, and
which finite family of tests decides between them. We fix all three.
Logical identity is theory-level Horn entailment on ordered premise
traces; behavioral identity is two-sided contextual indistinguishability
at the answer position; both are congruences on the same free monoid,
and the model respects logic exactly when the first refines the second.
We prove in Lean, without Mathlib, that this refinement is equivalent to
the existence of a unique canonical monoid morphism from the logical
quotient to the behavioral one; that premise reordering, repetition, and
derivable-clause extension are consequences of that single inclusion
rather than independent test relations; that a finite test family
decides the question exactly when the identity it induces is closed
under one-symbol extension, a closure failure naming the next test; and
that the observation table's biextensional collapse is the identifiable
object. We then build an instrument: frozen tables of prefix traces
against continuation--query tests, read through order relations on the
answer logits so that identity is exact, with preregistered statistics,
a non-reading validity control, and a surface-controlled contrast
between a maximal logic-preserving rewrite and a minimal logic-changing
one. On three base and instruction-tuned models up to 1.5B parameters,
no two logically identical traces are ever behaviorally identical;
models that read the task are moved as much by permuting or redundantly
extending the premises as by reversing one arrow. A recognizer built to
be permutation-invariant by architecture and trained to 99\% accuracy is
exactly the syntactic quotient and only partly the semantic one. Reading
the task, permutation invariance, and logical invariance are three
distinct properties, in that order of difficulty; neither accuracy nor
architectural invariance supplies the third.
\end{abstract}

\hypertarget{introduction}{%
\section{1. Introduction}\label{introduction}}

\textbf{The question.} Language models are sensitive to the order of
premises \cite{chen2024premise}, inconsistent across logically equivalent
reformulations \cite{zhou2026lgmt,gu2026crtbench}, and their accuracy and their
consistency come apart. These findings share a form: a set of inputs
that a logic declares identical, and a model that treats them
differently. What is missing is a statement of the object being tested.
Which identity on inputs? Which identity on outputs? And since only
finitely many inputs are ever tried, when has the test family decided
the question rather than sampled it?

\textbf{What we do.} We give the three definitions, prove what follows
from them, build an instrument that measures exactly what the
definitions name, and use it on measured and constructed recognizers.

\textbf{Contributions.} (i) A machine-checked framework in which
``respects logical identity'' is the refinement of one congruence by
another on a free monoid, with the equational consequences, the
finite-test identification criterion, and the collapse of the
observation table all proved (\texttt{recognition-paths}, Lean 4 core,
no Mathlib). (ii) A readout of answer logits through order relations
only, threshold-free and offset-invariant, which makes identity exact
and immune to answer bias, together with counts of exactly identical
profiles and collapse sizes. (iii) A surface-controlled statistic,
permutation against single-arrow flip, with a non-reading validity
control and a synthetic calibration. (iv) Constructed recognizers that
realise the syntactic half of logical identity by architecture, on which
the semantic half is measured in isolation.

\textbf{What we do not claim.} Order sensitivity and
accuracy--consistency gaps are known; our small-model findings replicate
them. The mathematics is classical (Myhill--Nerode, active automata
learning, and Chu spaces) \cite{angluin1987learning,pratt1995chu}; we assemble
it and check it. Nothing here concerns frontier
models.

\hypertarget{framework-proved}{%
\section{2. Framework (PROVED)}\label{framework-proved}}

\hypertarget{traces-queries-logical-identity}{%
\subsection{Traces, queries, logical
identity}\label{traces-queries-logical-identity}}

Fix a type of atoms. A Horn clause is a pair of atom lists (body, head);
a trace is a list of clauses; a query is a list of hypothesis atoms and
a goal atom. The theory of a trace, Γ(w), is its set of clauses.
Entailment is semantic: every valuation modelling Γ(w) and the
hypotheses makes the goal true. Two traces are logically identical, u
≡\_L v, when they entail the same queries. This is an equivalence and a
two-sided congruence for concatenation; since Γ forgets order and
repetition it is commutative and idempotent, so L = Σ*/≡\_L is a
commutative idempotent monoid. Permutations are logically identical;
appending a derivable clause is logically invisible
(\texttt{Horn.lean}).

\hypertarget{recognizers-and-behavioral-identity}{%
\subsection{Recognizers and behavioral
identity}\label{recognizers-and-behavioral-identity}}

A recognizer is ρ : Σ* × Q → O. Right-context identity u ≡\_ρ v asks
that no continuation and query separate u from v; two-sided identity u
≈\_ρ v asks the same in every left and right context. The latter is a
congruence, so B = Σ*/≈\_ρ is a monoid. The Nerode quotient Σ*/≡\_ρ is
the extensional collapse of the prefix realization (states = traces,
tests = continuation-- query pairs), and B acts on it as its transition
monoid; every state is reached from the empty trace
(\texttt{Recognition.lean}, \texttt{Nerode.lean}).

\hypertarget{the-factorization-theorem-and-its-consequences}{%
\subsection{The factorization theorem and its
consequences}\label{the-factorization-theorem-and-its-consequences}}

≡\_L ⊆ ≈\_ρ holds iff there is a unique representative-preserving F : L
→ B; F is then a monoid morphism. The reverse inclusion gives G : B → L;
both give L ≃ B. Under ≡\_L ⊆ ≈\_ρ, swapping two blocks of premises,
repeating a block, permuting a block, and appending a derivable clause
are all behaviorally invisible in every context. In the vocabulary of
metamorphic testing: reordering, duplication, and redundancy relations
are not independent; they are theorems of one inclusion
(\texttt{Factorization.lean}, \texttt{Recognition.lean}).

\hypertarget{when-a-finite-test-family-has-decided}{%
\subsection{When a finite test family has
decided}\label{when-a-finite-test-family-has-decided}}

Let $T$ be a family of tests $(z,q)$ containing the direct queries
$([],q)$. The identity it induces, $\equiv_{\rho,T}$, equals
$\equiv_\rho$ if and only if it is closed: $u\equiv_{\rho,T}v$ implies
$ua\equiv_{\rho,T}va$ for every symbol $a$. When $T$ is not closed there
exist $u\equiv_{\rho,T}v$ and a test $(a\cdot z,q)$ with $(z,q)\in T$
separating them, so the criterion is constructive. For
length-bounded families, one more level of continuation is one more
symbol of prefix, and a stable level identifies ≡\_ρ. Two-sided versions
hold for ≈\_ρ (\texttt{Identification.lean}). The counting bound (n
Nerode classes force stabilisation by length n−1) is standard and not
formalised (OPEN).

\hypertarget{the-table-as-a-chu-space}{%
\subsection{The table as a Chu
space}\label{the-table-as-a-chu-space}}

Quotienting tests by equal state profiles, dually to the state collapse,
gives the biextensional collapse, in which tests separate states and
states separate tests; silent extensions do not change it
(\texttt{Biextensional.lean}). The counts ``distinct rows'' and
``distinct columns'' of an observation table are the sizes of these
quotients. Under the broadest reading of ``geometry'', this point--test
duality is the geometry that cannot be removed; metric structure is a
further choice.

\hypertarget{specifications-for-constructed-recognizers}{%
\subsection{Specifications for constructed
recognizers}\label{specifications-for-constructed-recognizers}}

A theory-factoring recognizer ρ(w, q) = f(Γ(w), q) is invariant under
permutation and repetition of blocks by construction; the ideal
recognizer ρ(w, q) = Entails(Γ(w), q) is logically invariant and
recoverable, so L ≃ B. The gap between them, identification of
consequence-equivalent but syntactically different clause sets, is what
training must supply (\texttt{Specification.lean}).

\hypertarget{instrument}{%
\section{3. Instrument}\label{instrument}}

\hypertarget{tables}{%
\subsection{Tables}\label{tables}}

Eight Horn logical classes over five atoms, chosen so that the
class-level gold matrix has full rank and half of its cells are
positive. Rows are prefix traces (serializations, and in later tables
doubled prefixes, single-arrow flips, derivable-clause extensions);
columns are continuation--query pairs. Four frozen tables
(\texttt{hankel\_v0}--\texttt{v3}), each with a generator, an
independent validator that re-derives every gold label, a SHA-256 lock,
and CI. Row equivalences inside a class are Lean-certified; cell gold
labels are forward-chaining, not Lean-certified per cell.

\hypertarget{readout}{%
\subsection{Readout}\label{readout}}

Each cell records the logit pair (positive candidate, negative
candidate) at the answer position, one forward pass, float32, no
sampling. The analysis reads only order relations: the decision {[}pos
\textgreater{} neg{]} and, for two queries after the same continuation,
the preference {[}margin(q1) \textgreater{} margin(q2){]}. This is
threshold-free and invariant to the common logit offset, which had made
the raw metric table rank one, and it is immune to the answer-bias
saturation that discrete readouts inherit. Distances are Hamming counts
of separating tests: the number of columns of the Chu space that
distinguish two rows.

\hypertarget{statistics}{%
\subsection{Statistics}\label{statistics}}

All preregistered before the corresponding runs, with predictions
recorded. Gate R: comparative accuracy on query pairs whose gold labels
differ. Gate I: within-class median Hamming over between-class median.
S: median over classes of d\_perm / d\_flip on the surface-controlled
table. T, U, V: on the semantic-rewrite table, derivable extension
against flip, length-matched swap against flip, and the same on columns
where gold ties so that accuracy imposes nothing. E: counts of exactly
identical profiles. A non-reading recognizer (Pythia-70M) is the
validity control; a synthetic gold-plus-noise recognizer is the
calibration showing that T and U are passed by accuracy alone while V is
not.

\hypertarget{three-corrections}{%
\subsection{Three corrections}\label{three-corrections}}

The instrument corrected itself three times, each recorded. A Euclidean
readout was dominated by a logit offset (rank one) and by scale
differences across renderers; the Boolean readout replaced it. Gate I,
valid against the control, was found to reward shared surface content:
the class pair differing only in the direction of all three arrows was
closer than a class to its own serializations; the surface-controlled
contrast replaced it as the primary statistic. The apparent invisibility
of repeated continuations was found in the control at the same rate and
downgraded to a property of the prompts.

\hypertarget{measured-recognizers-observed}{%
\section{4. Measured recognizers
(OBSERVED)}\label{measured-recognizers-observed}}

Recognizers: Pythia-70M and Pythia-410M (base, \texttt{step143000}),
Qwen2.5-0.5B-Instruct and Qwen2.5-1.5B-Instruct (raw prompts, no chat
template), one forward pass per prompt in float32 on CPU, both answer
logits recorded.

\hypertarget{reading-is-not-identifying}{%
\subsection{Reading is not
identifying}\label{reading-is-not-identifying}}

On the Boolean table (\texttt{hankel\_v1}), Gate R and Gate I with their
controls:

\begin{longtable}[]{@{}
  >{\raggedright\arraybackslash}p{(\columnwidth - 4\tabcolsep) * \real{0.2727}}
  >{\raggedleft\arraybackslash}p{(\columnwidth - 4\tabcolsep) * \real{0.3636}}
  >{\raggedleft\arraybackslash}p{(\columnwidth - 4\tabcolsep) * \real{0.3636}}@{}}
\toprule\noalign{}
\begin{minipage}[b]{\linewidth}\raggedright
recognizer
\end{minipage} & \begin{minipage}[b]{\linewidth}\raggedleft
comparative acc. (Gate R \textgreater{} 0.55)
\end{minipage} & \begin{minipage}[b]{\linewidth}\raggedleft
within/between median Hamming (Gate I \textless{} 1)
\end{minipage} \\
\midrule\noalign{}
\endhead
\bottomrule\noalign{}
\endlastfoot
Pythia-70M & 0.456 & 1.42 \\
Qwen2.5-0.5B & 0.867 & 0.561 \\
\end{longtable}

Qwen reads the task and passes Gate I; the control fails both, so the
gate discriminates. But no recognizer has a single logical class whose
serializations have identical profiles (0 of 8 for every recognizer, all
40 rows distinct), and under a matched between-class statistic the
control also falls below 1. Gate I rewards shared surface content: for
Qwen the class pair differing only in the direction of all three arrows
is closer than a class to its own serializations. This is CRTBench's
accuracy-- consistency gap at small scale, with the addition that the
readout is exact: identity never occurs.

\hypertarget{surface-outweighs-logic-and-scale-moves-it}{%
\subsection{Surface outweighs logic, and scale moves
it}\label{surface-outweighs-logic-and-scale-moves-it}}

On the surface-controlled table (\texttt{hankel\_v2}), \texttt{S} =
median over classes of d\_perm / d\_flip:

\begin{longtable}[]{@{}lrrrr@{}}
\toprule\noalign{}
recognizer & \texttt{S} pooled & bullets & prose & comparative acc. \\
\midrule\noalign{}
\endhead
\bottomrule\noalign{}
\endlastfoot
Pythia-70M (control) & 0.97 & 1.20 & 1.11 & 0.48 \\
Qwen2.5-0.5B & 1.27 & 0.75 & 1.80 & 0.84 \\
Qwen2.5-1.5B & 0.99 & 0.80 & 1.25 & 0.89 \\
\end{longtable}

The statistic is one-sided: an accurate recognizer with no invariance
reaches 0.5--0.6 because flips change gold decisions it tracks. So the
diagnostic values are the ones above 1: for the 0.5B reader, and under
the prose renderer at both scales, reordering the same three premises
moves behavior more than reversing one arrow. From 0.5B to 1.5B the
pooled value falls from 1.27 to 0.99 (HYPOTHESIS: logic rises over
surface with scale within the family; two points).

\hypertarget{semantic-rewrites-are-not-identified}{%
\subsection{Semantic rewrites are not
identified}\label{semantic-rewrites-are-not-identified}}

On the semantic-rewrite table (\texttt{hankel\_v3}), a derivable-clause
extension (Lean-certified logically invisible) against a flip:

\begin{longtable}[]{@{}lrrr@{}}
\toprule\noalign{}
recognizer & \texttt{T} red/flip & \texttt{V} free columns & \texttt{E}
exact identities \\
\midrule\noalign{}
\endhead
\bottomrule\noalign{}
\endlastfoot
Pythia-70M (control) & 1.00 & 1.50 & 0 / 14 \\
Qwen2.5-0.5B & 1.11 & 1.50 & 0 / 14 \\
Qwen2.5-1.5B & 0.62 & 1.01 (bullets 0.67, prose 1.10) & 0 / 14 \\
\end{longtable}

For the 0.5B reader a logically invisible extension moves behavior as
much as a logical change, and more where accuracy imposes nothing
(\texttt{V} 1.5, equal to the control). At 1.5B \texttt{V} falls to the
no-identification value pooled and to 0.67 under the list renderer, with
no exact identification (HYPOTHESIS: identification on accuracy-free
columns begins with scale, renderer first). The behavioral quotient of
the 0.5B model does not contain the consequence relation at any level
the table sees.

\hypertarget{closure-idempotence-and-what-is-generic}{%
\subsection{Closure, idempotence, and what is
generic}\label{closure-idempotence-and-what-is-generic}}

The depth-one test family is not closed for any measured recognizer; for
Qwen the witness pair is logically identical and separated only at depth
two, the situation the identification theorem describes. Repeated
continuations agree with their single versions on about 95\% of tests
for reader and control alike, so this is a property of the prompts, not
of logic. The largest single source of within-class distance is a
first-premise primacy effect on the query mentioned first (CONFOUND),
and the atom-label family changes the metric invariance defect by about
3× (CONFOUND). Premise-by-premise margin increments are right-skewed,
heavy-tailed, drift toward the positive answer, do not accumulate with
depth, and shrink with level: not a diffusion (EXPLORATORY).

\hypertarget{constructed-recognizers}{%
\section{5. Constructed recognizers}\label{constructed-recognizers}}

Three recognizers trained on synthetic Horn traces with the eight
evaluation classes held out up to atom relabelling, 6000 steps, batch
128: \texttt{set} (max-pooled clause encodings: permutation- and
repetition-invariant by construction, \textasciitilde170k parameters),
\texttt{seq\_aug} (causal transformer with permutation/repetition
augmentation, \textasciitilde100k) and \texttt{seq\_fixed} (the same
without augmentation). Evaluated on the same tables through the same
statistics; one rendering, so no renderer split.

\hypertarget{the-syntactic-half-exactly}{%
\subsection{The syntactic half,
exactly}\label{the-syntactic-half-exactly}}

The \texttt{set} recognizer is, on \texttt{hankel\_v1}, exactly the
syntactic quotient: its 40 rows collapse to the 8 logical classes, all
eight classes have identical serialization profiles, \texttt{S} = 0 on
\texttt{hankel\_v2}, and the depth-one family is closed. Both seeds. No
causal recognizer, at any of four seeds or at three times the budget
(validation accuracy 0.98), has one class with identical profiles or
fewer than 39 distinct rows. Augmentation buys approximate permutation
invariance (\texttt{S} 0.49 at convergence); architecture buys exact
invariance; accuracy buys neither.

\hypertarget{the-semantic-half-partly}{%
\subsection{The semantic half,
partly}\label{the-semantic-half-partly}}

On \texttt{hankel\_v3} the \texttt{set} recognizer, 99\% accurate on
classes it never saw, identifies none of the 14 derivable-clause
extensions exactly in one seed and two in the other; on the columns
where accuracy imposes nothing it moves 0.62--0.80 times as far under a
semantic rewrite as under a logical flip. Every LLM, the control, and
seven of eight causal toy runs have \texttt{V} above 1 (medians
1.5--2.3). So the architecturally invariant recognizer is the only one
in which a semantic rewrite is detectably less disruptive than a logical
change, and even there most of the semantic half is missing and almost
none of it is exact. (Two seeds; suggestive.)

\hypertarget{what-the-constructed-recognizers-add}{%
\subsection{What the constructed recognizers
add}\label{what-the-constructed-recognizers-add}}

They separate three properties that the LLM measurements conflate:
reading the task (all accurate recognizers), permutation invariance
(exact only by architecture, approximate by augmentation), and logical
invariance (present partly, and only where the syntactic half was
exact). They also show that \texttt{S\ \textless{}\ 1} and Gate I are
reached by accuracy and surface alone, which fixes the reading of
Section 4.

\hypertarget{discussion}{%
\section{6. Discussion}\label{discussion}}

\textbf{Three properties, ordered.} Reading the task, respecting the
syntactic part of logical identity, and respecting its semantic part are
distinct, and every recognizer we measured or built sits on that ladder
at a different rung. Accuracy does not climb it; architecture climbs one
rung exactly; nothing we tried climbs the last.

\textbf{What a finite table can conclude.} By the identification
theorem, a test family decides \texttt{≡\_ρ} when it is closed. On the
measured recognizers the depth-one family is not closed, and the theorem
names the column a deeper table needs. On the \texttt{set} recognizer
the same family is closed, which is the theorem working as intended: the
recognizer's quotient is coarse enough for the family to capture it.

\textbf{Exact versus approximate.} A statistical recognizer never
produces two identical profiles for logically identical inputs; the
exceptions here are forced by construction. The right long-run object is
a canonical distance with a certified tolerance, not an equality (OPEN).

\textbf{Metric and duality.} Metric geometry on the logits misled twice
(offset, scale); the Boolean readout keeps only the point--test duality
of the Chu space, and every exact statement in this paper lives there.

\textbf{Limitations.} One Horn fragment, five atoms, eight classes;
models up to 1.5B without chat templates; no sampling; toy constructed
models at one or four seeds; gold labels by forward chaining, not
certified per cell.

\textbf{Related work and novelty.} Prior work establishes premise-order
sensitivity \cite{chen2024premise}, logic-grounded metamorphic testing
\cite{zhou2026lgmt}, controlled reformulation consistency
\cite{gu2026crtbench}, parameterized logical robustness
\cite{essebbani2026robustness}, permutation-invariant architectures
\cite{egressy2025setllm}, and order-centric augmentation
\cite{he2025order}. Weighted-automata extraction also supplies a nearby
Hankel-matrix tradition \cite{lacroce2021extracting}. The empirical phenomena
are therefore known. Our narrower contribution is the machine-checked
quotient framework, the finite-test closure criterion, the exact ordinal
readout, the controlled surface statistic, and the experimental separation
of permutation invariance from logical invariance.

\section{Artifacts and reproducibility}
The Lean 4 formalization is archived at
\url{https://doi.org/10.5281/zenodo.22495808}. Frozen benchmark generators,
validators, lock files, model manifests, raw measurements, analysis scripts,
and constructed-recognizer checkpoints are archived at
\url{https://doi.org/10.5281/zenodo.22495706} and mirrored at
\url{https://github.com/Kairose-master/proof-path-invariance}. The repository's
continuous-integration workflow rebuilds the formal certificates, regenerates
the benchmark families, independently validates their labels, and verifies
their recorded hashes.
